\clearpage
\phantomsection% 
\addcontentsline{toc}{chapter}{ABSTRAK}
\thispagestyle{fancy}

\begin{center}
	\large \bfseries \MakeUppercase{Abstrak}\\
	\normalsize \normalfont {\thetitle}\\
	\normalsize \normalfont {\theauthor}\\
	\bigskip
	
	\normalsize \normalfont \justifying \singlespacing
	Kemajuan pengamatan astronomi telah beralih ke sistem otomatisasi, namun banyak teleskop otomatis menghadapi tantangan pada konfigurasi awal dan kalibrasi, terutama bagi pengguna dengan keterampilan teknis terbatas. Penelitian ini mengembangkan teleskop robotik berbasis dudukan Alt-Azimuth dengan sistem kontrol otomatis melalui aplikasi Android sebagai antarmuka utama. Sistem mengintegrasikan sensor IMU MPU6050, magnetometer QMC5883L, dan encoder magnetik AS5600 untuk memberikan informasi orientasi secara real-time, memungkinkan teleskop menyesuaikan posisi secara otomatis saat dinyalakan atau mengalami pergeseran. Untuk mengatasi keterbatasan sensor MPU6050 (noise, drift, dan ketidakakuratan), diimplementasikan Mahony Filter yang menggabungkan data giroskop dan akselerometer menggunakan pendekatan pengendali PI (Proporsional-Integral) untuk menghasilkan estimasi orientasi yang lebih akurat. Kinerja sistem dievaluasi menggunakan metrik Mean Absolute Error (MAE) dan Root Mean Square Error (RMSE). Hasil menunjukkan Mahony Filter berhasil mereduksi noise Roll sebesar 79,1\% (dari 0,510$^{\circ}$ menjadi 0,107$^{\circ}$) dan Pitch sebesar 84,2\% (dari 1,719$^{\circ}$ menjadi 0,271$^{\circ}$). Sistem GoTo berhasil mengarahkan teleskop ke objek langit dengan pointing error 2,52$^{\circ}$, sedangkan sistem tracking dapat mengikuti pergerakan semu benda langit selama 1 jam, meskipun belum dapat mempertahankan objek secara konsisten di pusat bidang pandang akibat faktor mekanik (flex struktur PLA dan backlash gear). Sistem ini memadai untuk aplikasi astronomi amatir dan edukasi dengan objek langit terang berdiameter sudut besar (Bulan, Matahari), namun masih memerlukan perbaikan mekanik signifikan untuk pengamatan objek titik (planet, bintang) yang memerlukan presisi lebih tinggi.\par

    \vspace{20pt}
	\raggedright \textbf{Kata Kunci: Teleskop Robotik, Dudukan Alt-Azimuth, Mahony Filter, Sensor IMU, Sistem GoTo}
	
	\vfill
	
\end{center}
\clearpage